% 小行星（综述）
% license CCBYSA3
% type Wiki

本文根据 CC-BY-SA 协议转载翻译自维基百科\href{https://en.wikipedia.org/wiki/Asteroid}{相关文章}。

\begin{figure}[ht]
\centering
\includegraphics[width=8cm]{./figures/9fdf35f418301f87.png}
\caption{被探测访问过的小行星影像展示其差异：（上排）433 Eros 与 243 Ida 及其卫星 Dactyl，（下排）Ceres 与 101955 Bennu。尺寸不按比例绘制。} \label{fig_Astero_1}
\end{figure}
小行星是一类小行星体（minor planet），其尺寸大于流星体（meteoroid），但既不是行星，也不是已确认为彗星的天体。它们位于太阳系内部轨道运行，或与木星共享轨道（如特洛伊小行星）。小行星由岩石、金属或冰质物质构成，不具备大气层，并可大致分类为 C 型（碳质）、M 型（金属）、S 型（硅酸盐）。小行星的大小与形状差异极大，从直径不足一千米的松散碎石堆，到直径接近 1000 千米、被归类为矮行星的 Ceres。若天体在受到太阳辐射加热时产生彗发（尾），则被归类为彗星而非小行星，尽管近期观测显示这两类天体之间可能存在连续光谱。\(^\text{[1][2]}\)

在约一百万颗已知小行星中，\(^\text{[3]}\)数量最多的位于火星与木星轨道之间、距离太阳约 2 至 4 AU 的区域，称为主小行星带（main asteroid belt）。所有小行星总质量仅约为地球月球质量的 3\%。主带中大多数小行星沿轻微椭圆、稳定的轨道运行，方向与地球一致，绕太阳一周约需 3 至 6 年。\(^\text{[4]}\)

小行星最早通过地面观测方式进行研究，首次近距离观测来自伽利略号（Galileo）探测器。其后，美国 NASA 与日本 JAXA 相继发射了多次专门的小行星探测任务，并持续规划新的任务。NASA 的 NEAR Shoemaker 研究了 Eros，Dawn 探测了 Vesta 与 Ceres；JAXA 的隼鸟号（Hayabusa）与隼鸟 2 号分别研究并回收了 Itokawa 与 Ryugu 的样本；OSIRIS{-}REx 探测 Bennu，并于 2020 年采样，2023 年将样本送回地球。NASA 的 Lucy（2021 年发射）将研究十颗不同小行星，其中两颗来自主带、八颗为木星特洛伊小行星。Psyche 探测器于 2023 年 10 月发射，目标是金属小行星 Psyche。欧洲航天局的 Hera 于 2024 年 10 月发射，旨在研究 DART 碰撞后的结果。中国国家航天局的天问二号于 2025 年 5 月发射，\(^\text{[5]}\)其任务为探测共轨近地小行星 469219 Kamoʻoalewa 与活跃小行星 311P/PanSTARRS，并对 Kamoʻoalewa 的风化层（regolith）进行采样。\(^\text{[6]}\)

近地小行星若与地球发生撞击，可能造成灾难性影响，其中最著名的案例是 Chicxulub 撞击事件，被普遍认为导致了白垩纪–古近纪生物大灭绝。为应对这一潜在威胁，作为实验性防御措施，2022 年 9 月，双小行星偏转测试（Double Asteroid Redirection Test, DART）探测器通过撞击成功改变了无威胁小行星 Dimorphos 的轨道。
\subsection{术语}
\begin{figure}[ht]
\centering
\includegraphics[width=8cm]{./figures/795a52138582bb97.png}
\caption{一幅按相同尺度制作的组合图像，展示了 2012 年前以高分辨率成像的小行星。从大到小依次为：4 Vesta、21 Lutetia、253 Mathilde、243 Ida 及其卫星 Dactyl、433 Eros、951 Gaspra、2867 Šteins、25143 Itokawa。} \label{fig_Astero_2}
\end{figure}
\begin{figure}[ht]
\centering
\includegraphics[width=8cm]{./figures/1f1fc84d80338171.png}
\caption{按相同比例展示的 Vesta（左）、Ceres（中）与月球（右）} \label{fig_Astero_3}
\end{figure}
2006 年，国际天文学联合会（IAU）引入了当前更为推荐的宽泛术语“太阳系小天体”（small Solar System body），将其定义为太阳系中既不是行星、也不是矮行星、亦不是天然卫星的天体；其中包括小行星、彗星以及更近期发现的其他类型天体\(^\text{[7]}\)。根据 IAU 表述，“术语『小行星（minor planet）』仍可继续使用，但一般情况下将优先采用『太阳系小天体（Small Solar System Body）』”\(^\text{[8]}\)。

在历史上，第一颗被发现的小行星——Ceres——最初曾被视为一颗新行星\(^\text{[a]}\)。随后，天文学家又陆续发现了其他相似天体。在当时的观测设备条件下，它们呈现为类似恒星的光点，几乎看不到行星盘结构，但由于具有明显的视运动，又可与恒星区分开来。这促使天文学家威廉·赫歇尔（Sir William Herschel）提出了“asteroid”一词\(^\text{[b]}\)，该词源自希腊语 ἀστεροειδής（asteroeidēs），意为“类星”“似星状”，并由古希腊语 ἀστήρ（astēr，“星、行星”）派生。在 19 世纪后半叶初期，“asteroid”与“planet”（有时并未特别限定为“minor”）仍被交替使用\(^\text{[c]}\)。

传统上，围绕太阳运行的小型天体被划分为彗星、小行星或流星体，其中直径小于一米的天体被称为流星体。术语“asteroid”从未被正式定义\(^\text{[13]}\)，但在非正式使用中可指“围绕太阳运行、形状不规则、呈岩质组成且不符合 IAU 定义中行星或矮行星条件的天体”\(^\text{[14]}\)。小行星与彗星的主要区别在于，彗星因近表层冰在太阳辐射作用下发生升华而形成彗发（尾）。少数天体最初被归类为小行星，但后续观测显示存在彗状活动。相反，一些（或许全部）彗星最终会耗尽其表面挥发性冰，从而演化为类小行星天体。另一区分方式是轨道偏心率：彗星通常具有比大多数小行星更高的轨道偏心率；而高度偏心的小行星极可能为休眠或灭绝的彗星\(^\text{[15]}\)。

位于木星轨道之外的小行星体有时亦被称为“asteroids”，尤其是在大众化科普材料中\(^\text{[d]}\)。然而，将“asteroid”一词限定用于太阳系内部的小行星体的做法正在变得愈发普遍\(^\text{[17]}\)。因此，本文主要将讨论范围限定于传统意义上的小行星：即小行星带天体、木星特洛伊小行星以及近地天体。

自 1801 年 Ceres 被发现后的近两个世纪中，所有已知小行星在其轨道的大部分时间里都位于木星轨道以内，尽管极少数如 944 Hidalgo 会在轨道部分区段上延伸至更远处。自 1977 年 2060 Chiron 被发现起，天文学家开始发现永久位于木星轨道之外的小型天体，即后称为半人马（centaurs）。1992 年，又发现了 15760 Albion，这是首个已知位于海王星轨道之外的天体（除 Pluto 外）；随后人们观测到了大量类似天体，即现称为跨海王星天体（trans-Neptunian object）。更外层依次还有柯伊伯带天体、散射盘天体以及距离更为遥远的奥尔特云——其被假设为休眠彗星的主要储库。它们位于太阳系寒冷的外缘区域，冰质物质能够长期保持固态，因此彗状活动极为微弱；若半人马或跨海王星天体接近太阳，其挥发性冰质材料将发生升华，按传统分类方法将被归类为彗星。

之所以将柯伊伯带天体称为“objects”，部分原因是为了避免必须将其归为小行星或彗星\(^\text{[1]}\)。它们被认为在成分上大体类似于彗星，但部分天体可能与小行星更相近\(^\text{[18]}\)。绝大多数此类天体并不具有典型彗星所表现的高偏心率轨道，且目前发现者均比传统彗核更大。近年来的观测结果，如 Stardust 探测器采集的彗星尘埃样本分析，正不断模糊小行星与彗星之间的界限\(^\text{[1]}\)，越来越多研究表明二者之间可能形成连续谱，而非清晰的分类边界\(^\text{[2]}\)。

2006 年，IAU 为最大的小行星体引入了“矮行星”这一类别——即因自身引力而形成近似椭球形状的天体。迄今为止，仅小行星带中体积最大的天体 Ceres（直径约 975 km，约 606 mi）被划入此类\(^\text{[19][20]}\)。
\subsection{观测历史}
尽管小行星数量极为庞大，但其发现在人类天文学史上相对较晚，第一颗小行星——Ceres——直到 1801 年才被识别出来\(^\text{[21]}\)。在暗夜且观测条件良好的情况下，只有一颗小行星，即表面反照率较高的 4 Vesta，通常能够以肉眼看见。极少数情况下，当小型小行星非常接近地球时，也可能在短暂时间内被肉眼观测到\(^\text{[22]}\)。截至 2025 年 5 月，小行星中心（Minor Planet Center）已记录太阳系内外共计 1,460,356 颗小行星的数据，其中约有 826,864 颗具备足够观测信息而获得编号编号\(^\text{[23]}\)。
\subsubsection{Ceres 的发现}
\begin{figure}[ht]
\centering
\includegraphics[width=6cm]{./figures/63dcf87fd16a81c7.png}
\caption{《Giuseppe Velasquez 与 Giuseppe Piazzi 及 Ceres》，油画，约作于 1803 年} \label{fig_Astero_4}
\end{figure}
1772 年，德国天文学家 Johann Elert Bode（援引 Johann Daniel Titius 的观点）发表了一组数列进程，即如今已被否定的 Titius–Bode 定律。除了火星与木星之间存在无法解释的空隙外，该公式似乎能够预测当时已知行星的轨道\(^\text{[24][25]}\)。Bode 对这一“缺失行星”的存在给出了如下解释：

“这一点尤其来源于一项令人惊异的关系，即已知六颗行星在其距日距离上所遵循的规律。若取土星到太阳的距离为 100，则水星距太阳为 4；金星为 4 + 3 = 7；地球为 4 + 6 = 10；火星为 4 + 12 = 16。此处出现一个在规则进程中的空档：在火星之后应为 4 + 24 = 28 的位置，但迄今未见任何行星。难道可以相信宇宙的创造者故意让这一区域保持空缺吗？显然不能。从此处继续，木星距离为 4 + 48 = 52，最终土星为 4 + 96 = 100。”\(^\text{[26]}\)

依据该公式，Bode 预测将在距离太阳约 2.8 天文单位（AU），即约 4.2 亿千米的位置发现另一颗行星\(^\text{[25]}\)。随着 William Herschel 在接近公式预测位置发现了土星外的新行星天王星，Titius–Bode 定律一度得到了强有力支持\(^\text{[24]}\)。1800 年，由德国天文学刊物《Monatliche Correspondenz》（《每月通信》）主编 Franz Xaver von Zach 牵头，他向 24 名资深天文学家发出邀请（并将其称为“天国警察”）\(^\text{[25]}\)，希望他们协同开展针对该“预期行星”的系统性搜寻\(^\text{[25]}\)。虽然他们未能发现 Ceres，但最终陆续发现了编号为 2、3、4 的三颗小行星：2 Pallas、3 Juno 与 4 Vesta\(^\text{[25]}\)。

被选入搜寻小组的天文学家之一是 Giuseppe Piazzi，他是西西里巴勒莫学院的一位天主教神父。在收到加入该小组的正式邀请之前，Piazzi 已于 1801 年 1 月 1 日发现了 Ceres\(^\text{[27]}\)。当时他正在寻找“拉卡伊先生（Mr la Caille）《黄道星目录》中第 87 颗星”\(^\text{[24]}\)，却注意到“其前方又出现了一颗天体”\(^\text{[24]}\)。Piazzi 实际上发现了一颗会移动的、类星状的天体，并最初认为其是一颗彗星\(^\text{[28]}\)：

“其光度略显微弱，颜色类似木星，与通常被视为八等星的许多恒星非常相近。因此，我毫不怀疑其应为一颗普通恒星。[……]然而在第三天夜晚，我的怀疑已转为确信——它绝非一颗固定恒星。不过在公开之前，我仍决定等待至第四天傍晚，当时我欣喜地看到它继续以前几日相同的速率移动。”\(^\text{[24]}\)

Piazzi 共计观测 Ceres 24 次，最后一次记录于 1801 年 2 月 11 日，之后因病中断研究。他于 1801 年 1 月 24 日向两位同行天文学家致信宣布此发现：米兰的同乡 Barnaba Oriani 与柏林的 Bode\(^\text{[21]}\)。在信中，他将其报告为一颗彗星，但同时表示：“鉴于其运动如此缓慢且相当规则，我已多次想到它可能比彗星更为特殊。”\(^\text{[24]}\)

4 月，Piazzi 将完整观测数据寄给 Oriani、Bode 与法国天文学家 Jérôme Lalande，并于当年 9 月发表于《Monatliche Correspondenz》期刊\(^\text{[28]}\)。

其时，由于地球绕太阳公转引起的视位置变化，Ceres 的位置已发生偏移，并过于靠近太阳炫光，其他天文学家无法立即证实 Piazzi 的观测结果。年末 Ceres 理论上应再次可见，但由于间隔时间长、轨道参数尚不精确，其位置难以预测。为重新寻获 Ceres，时年仅 24 岁的数学家 Carl Friedrich Gauss 开发出一种高效的轨道确定方法\(^\text{[28]}\)。仅数周后，他完成预测并将结果寄给 von Zach。1801 年 12 月 31 日，von Zach 与另一位“天国警察”Heinrich W. M. Olbers 在几乎完全符合预测的位置成功重新发现 Ceres\(^\text{[28]}\)。Ceres 位于距太阳 2.8 AU 的轨道位置，与 Titius–Bode 定律几乎完全吻合；然而海王星在 1846 年被发现后，其距离预测值偏差高达 8 AU，促使大多数天文学家最终将该定律视为巧合\(^\text{[29]}\)。Piazzi 将新发现的天体命名为 Ceres Ferdinandea，以“纪念西西里守护女神与波旁王朝国王 Ferdinand”\(^\text{[26]}\)。
\subsubsection{进一步搜寻}
\begin{figure}[ht]
\centering
\includegraphics[width=8cm]{./figures/321bbb2080ca89a0.png}
\caption{首十颗被发现的小行星的尺寸，与月球的比较} \label{fig_Astero_5}
\end{figure}

